\documentclass{uethesis}

\usepackage{booktabs}

%%%%%%%%%%%%%%%%%%%%%%%%%%%%%%%%%%%%%%%%%%%%%%%%%%%%%%%%%%%%%%%%%%%%
%% Set Thesis Information
%%%%%%%%%%%%%%%%%%%%%%%%%%%%%%%%%%%%%%%%%%%%%%%%%%%%%%%%%%%%%%%%%%%%
\renewcommand{\thesistitle}{Designing Lightweight Agentic AI for E-Commerce: From Intent Prediction to Marketing Optimization}
\renewcommand{\thesisauthor}{Pratik Chandrakant Gotakhinde}
\renewcommand{\matriculation}{88215836}
\renewcommand{\firstsupervisor}{Prof. Dr. Iftikhar Ahmad}
\renewcommand{\secondsupervisor}{Prof. Dr. Raja Hasim Ali}
\renewcommand{\submissiondate}{August xx, 2026}
\renewcommand{\department}{Department of Tech and Software Engineering}
\renewcommand{\major}{Data Science}
\renewcommand{\degree}{Master of Data Science} 

%%%%%%%%%%%%%%%%%%%%%%%%%%%%%%%%%%%%%%%%%%%%%%%%%%%%%%%%%%%%%%%%%%%%
%% Define Abstract Content Here
%%%%%%%%%%%%%%%%%%%%%%%%%%%%%%%%%%%%%%%%%%%%%%%%%%%%%%%%%%%%%%%%%%%%
\renewcommand{\theabstract}{
The abstract of a thesis should provide a concise summary of the research. It should begin by outlining the context and motivation, explaining the research problem or gap and its significance. Next, it should briefly describe the methodology, summarizing the approach or methods used to address the problem. The abstract should then highlight the key results and conclusions, emphasizing the main findings and their implications. Finally, it should mention the significance of the research, explaining its broader impact and contribution to the field.

The abstract should not contain any subsection, citations, or references to any text, figures, tables in the main body of thesis. It should be self contained. As a guideline the length should be at least 3/4th of a page, but should not exceed a single page.


  \vspace{0.5cm}
  \noindent \textbf{Keywords:} Lorem, Ipsum, Dolor, Sit, Amet
}

%%%%%%%%%%%%%%%%%%%%%%%%%%%%%%%%%%%%%%%%%%%%%%%%%%%%%%%%%%%%%%%%%%%%
%% Begin Document
%%%%%%%%%%%%%%%%%%%%%%%%%%%%%%%%%%%%%%%%%%%%%%%%%%%%%%%%%%%%%%%%%%%%
\begin{document}

% Begin Thesis Setup: Title page, front matter, etc.
\beginthesis

%%%%%%%%%%%%%%%%%%%%%%%%%%%%%%%%%%%%%%%%%%%%%%%%%%%%%%%%%%%%%%%%%%%%
%% Chapter 1: Introduction
%%%%%%%%%%%%%%%%%%%%%%%%%%%%%%%%%%%%%%%%%%%%%%%%%%%%%%%%%%%%%%%%%%%%
\chapter{Introduction}
The following text is only provided for information, change the sections, sub-sections and text as per the requirements of your thesis.

\section{}
LaTeX is a powerful typesetting system renowned for its ability to produce high-quality technical and scientific documents. In the following, we provide an example demonstrating how to insert and refer to a table (Table~\ref{tab:ml_results}) and an image (Fig~\ref{fig:confusion_matrix}) into your document. For more information on using LaTeX, please refer to the Overleaf24 tutorial~\citep{overleaf2024}.

% --- Table: Machine Learning Model Performance ---
\begin{table}[h!]
    \centering
    \caption{Performance Metrics for Various Machine Learning Models}
    \label{tab:ml_results}
    \begin{tabular}{lcccc}
        \toprule
        \textbf{Model} & \textbf{Accuracy} & \textbf{Precision} & \textbf{Recall} & \textbf{F1 Score} \\
        \midrule
        Random Forest           & 92.5\%  & 91.0\%  & 93.2\%  & 92.1\%  \\
        Support Vector Machine  & 90.3\%  & 89.5\%  & 91.0\%  & 90.2\%  \\
        Neural Network          & 93.7\%  & 92.8\%  & 94.1\%  & 93.4\%  \\
        \bottomrule
    \end{tabular}
\end{table}

% --- Figure: Confusion Matrix ---
\begin{figure}[h!]
    \centering
    \includegraphics[width=0.7\textwidth]{example-image} % Replace with your actual image file name
    \caption{Confusion Matrix for the Neural Network Model}
    \label{fig:confusion_matrix}
\end{figure}


\section{Background}
This section provides essential context and motivation for the research. 

\section{Problem Statement}
Here, we define the specific problem the thesis aims to address. 

\section{Research Objectives}
We outline the goals and objectives that guide the study. Sed hendrerit sapien 

\section{Research Questions}
We specify the key questions that the research seeks to answer. 

\section{Contributions}
A brief summary of the novel aspects or main contributions of the work. 

\section{Thesis Organization}
The rest of the thesis is organized as follows. In Chapter~\ref{chap:background}, the background of the study is presented, laying the groundwork by discussing the theoretical foundations and context in which the research is situated. Chapter~\ref{chap:lit} provides a comprehensive literature review, critically examining previous work and highlighting gaps that this research aims to address. In Chapter~\ref{chap:method}, the methodology is detailed, describing the research design, data collection, and analysis procedures employed in the study. Chapter~\ref{chap:res} presents the results and discussions, where the findings are analyzed in relation to the research objectives. Finally, Chapter~\ref{chap:conc} concludes the thesis by summarizing the key contributions, reflecting on the implications of the results, and suggesting directions for future research.


%%%%%%%%%%%%%%%%%%%%%%%%%%%%%%%%%%%%%%%%%%%%%%%%%%%%%%%%%%%%%%%%%%%%
%% Chapter 2: Background
%%%%%%%%%%%%%%%%%%%%%%%%%%%%%%%%%%%%%%%%%%%%%%%%%%%%%%%%%%%%%%%%%%%%
\chapter{Background} \label{chap:background}
\textcolor{red}{The following text is for guidance, replace it based on the requirements of your own work/thesis.}

This chapter lays the foundation for the study by introducing the key concepts, theories, and historical context relevant to the research area. It provides the necessary background information that situates the study within a broader academic and practical framework, ensuring that readers are well-prepared for the discussion that follows.

%%%%%%%%%%%%%%%%%%%%%%%%%%%%%%%%%%%%%%%%%%%%%%%%%%%%%%%%%%%%%%%%%%%%
%% Chapter 3: Literature Review
%%%%%%%%%%%%%%%%%%%%%%%%%%%%%%%%%%%%%%%%%%%%%%%%%%%%%%%%%%%%%%%%%%%%
\chapter{Literature Review}\label{chap:lit}
\textcolor{red}{The follwoing text is for guidance, replace it based on the requirements of your own work/thesis.}

In this chapter, a comprehensive review of the existing literature is presented. The discussion critically examines previous research and theoretical contributions, highlighting both established findings and identified gaps that motivate the current study. This analysis helps to position the research within the ongoing scholarly conversation.


%%%%%%%%%%%%%%%%%%%%%%%%%%%%%%%%%%%%%%%%%%%%%%%%%%%%%%%%%%%%%%%%%%%%
%% Chapter 4: Methodology
%%%%%%%%%%%%%%%%%%%%%%%%%%%%%%%%%%%%%%%%%%%%%%%%%%%%%%%%%%%%%%%%%%%%
\chapter{Methodology} \label{chap:method}
\textcolor{red}{The following text is for guidance, replace it based on the requirements of your own work/thesis.}

This chapter outlines the research design and the methods employed to gather and analyze data. It explains the rationale behind the chosen methodology, details the data collection procedures, and describes the analytical techniques used. The chapter aims to provide a clear and replicable roadmap for the study.



%%%%%%%%%%%%%%%%%%%%%%%%%%%%%%%%%%%%%%%%%%%%%%%%%%%%%%%%%%%%%%%%%%%%
%% Chapter 5: Results and Discussions
%%%%%%%%%%%%%%%%%%%%%%%%%%%%%%%%%%%%%%%%%%%%%%%%%%%%%%%%%%%%%%%%%%%%
\chapter{Results and Discussions} \label{chap:res}
\textcolor{red}{The following text is for guidance, replace it based on the requirements of your own work/thesis.}

Here, the findings of the research are systematically presented and interpreted. The chapter begins by outlining the key results, supplemented by relevant tables and figures, and then discusses these outcomes in light of the research questions and theoretical framework. This analysis bridges the gap between the raw data and the broader implications of the study.



%%%%%%%%%%%%%%%%%%%%%%%%%%%%%%%%%%%%%%%%%%%%%%%%%%%%%%%%%%%%%%%%%%%%
%% Chapter 6: Conclusions
%%%%%%%%%%%%%%%%%%%%%%%%%%%%%%%%%%%%%%%%%%%%%%%%%%%%%%%%%%%%%%%%%%%%
\chapter{Conclusions} \label{chap:conc}
\textcolor{red}{The following text is for guidance, replace it based on the requirements of your own work/thesis.}

This final chapter summarizes the main findings and reflects on the overall contributions of the research. It discusses the implications of the results, acknowledges any limitations, and offers recommendations for future research. The chapter serves to encapsulate the study’s significance and suggest practical steps forward.

\section{Conclusions}

\section{Future Works}

%%%%%%%%%%%%%%%%%%%%%%%%%%%%%%%%%%%%%%%%%%%%%%%%%%%%%%%%%%%%%%%%%%%%
%% Bibliography
%%%%%%%%%%%%%%%%%%%%%%%%%%%%%%%%%%%%%%%%%%%%%%%%%%%%%%%%%%%%%%%%%%%%
\clearpage
\bibliography{ref} % Provide ref.bib in the same folder

\end{document}
